\documentclass[10pt]{article}

\usepackage{parskip}
\usepackage[margin=1in]{geometry} 
\usepackage{amsmath,amsthm,amssymb, graphicx, multicol, array}
\usepackage{enumitem}
\usepackage{amssymb}
\usepackage{float}
\usepackage{href-ul}
 
\newcommand{\N}{\mathbb{N}}
\newcommand{\Z}{\mathbb{Z}}
 
\newenvironment{problem}[2][Problem]{\begin{trivlist}
\item[\hskip \labelsep {\bfseries #1}\hskip \labelsep {\bfseries #2.}]}{\end{trivlist}}

\begin{document}
 
\title{Introduction}
\author{Jakob Sverre Alexandersen\\
GRA4154 Supply Chain Optimization with Mathematical Programming\\
Lecture 1}
\maketitle

\tableofcontents
\newpage


\section{Introduction}
\subsection{Mathematical Programming}
\begin{itemize}
    \item A way to use math to make the best possible decision when there are rules to follow (constraints)
    \item You have choices you can control (like how much of each product to make)
    \item You have a goal (like maximizing profits or minimizing cost)
    \item You have constraints (limits such as time, money or resources)
    \item Math helps you find the best solution that satisfies all the limits 
\end{itemize}
\subsection{Production planning model for Orkla}

\begin{gather*}
    \text{Total cost} = \sum_{t = 1}^{W}\sum_{i = 1}^{N}\sum_{j = 1}^{N}C_{\text{real}}\times T_{ijt} + \sum_{t = 1}^{W}\sum_{i = 1}^{N}h_{\text{real}}\times I_{\text{real}}\\
    = \sum_{t = 1}^{W}\sum_{i = 1}^{N}\sum_{j = 1}^{N}s_{\text{real}}\times\text{unit labor cost}\times \text{\# of operators}\times T_{ijt}\\
    + \sum_{t = 1}^{W}\sum_{i = 1}^{N}h_{\text{real}}\times I_{\text{model}}\times\text{capacity / machine time }\\
    = \sum_{t = 1}^{W}\sum_{i = 1}^{N}\sum_{j = 1}^{N}s_{\text{model}}\times\text{capacity}\times\text{unit labor cost} \times\text{\# of operators}\times T_{ijt}\\
    + \sum_{t = 1}^{W}\sum_{i = 1}^{N} (h_{\text{real}})\times\text{capacity / machine time}\times I_{\text{model}}\\
    = \sum_{t = 1}^{W}\sum_{i = 1}^{N}\sum_{j = 1}^{N}C_{\text{model}}\times T_{ijt} + \sum_{t = 1}^{W}\sum_{i = 1}^{N}h_{\text{model}}\times I_{\text{model}}\\
    \text{s.t.}\\
\end{gather*}
\begin{gather}
    I_{i, t - 1} + X_{it} - I_{it} = D_{it}\quad\forall i, t\quad\text{(to balance the inventory)}\\
    X_{it} - Y_{it} \le 0\quad\forall i, t\quad\text{(to ensure whenever $X_{it} > 0 $ then $Y_{it} = 1$)}\\
    \sum_{i = 1}^{N}X_{it} + \sum_{i = 1}^{N}\sum_{j = 1}^{N}T_{ijt}s_{ij}\le 1\quad\forall t\quad\text{(capacity constraint)}
\end{gather}
% Constraint (X) to (Y) determine the first and last product to be produced in a period, and also find the product for which the machine is setup at the end of each period 
\begin{gather}
    Y_{it} \le\omega_t\quad\forall i, t\quad\text{($\omega_t = 1 $ if any item is produced in period $t $)}\\
    \sum_{i = 1}^{N}Y_{it} - 1\le (N - 1)\delta_t\quad\forall t\\
    \omega_t\le\sum_{i = 1}^{N}\alpha_t\le 1\quad\forall t\\
    \omega_t\le\sum_{i = 1}^{N}\beta_t\le 1\quad\forall t\\
    \alpha_{it}\le Y_{it}\quad\forall i, t\\
    \beta_{it}\le Y_{it}\quad\forall i, t
\end{gather}
\newpage 
\begin{itemize}
    \item Constraint (4) to (10) determine the first and last product to be produced in a period, and also find the product for which the machine is setup at the end of each period 
    \item (6) and (7) Ensure that $\alpha \land\beta = 1$ for exactly one product in a given period if production occurs in that period 
    \item (8) and (9) ensure that in a no-production perdiod, all $\alpha = \beta = 0 $
\end{itemize}
\begin{gather}
    \sum_{j = 1}^{N}\gamma_{it} = 1\quad\forall t\quad\text{(only one product can be set up at the end of each period)}\\
    \sum_{j = 1}^{N}T_{ijt}\ge Y_{it} - \alpha_{it}\quad\forall i, t\\
    \sum_{j = 1}^{N}T_{ijt}\ge Y_{it} - \beta_{it}\quad\forall i, t\\
    T_{ijt}\ge \gamma_{i, t - 1} - \omega_t - 1\quad\forall\neq j, t
\end{gather}
\begin{itemize}
    \item (11) and (12): Apply when $\ge 2 $ products are produced in a period: force at least one $T_{ijt} $ to be 1 for each product $i $, except when this product is either the first or the last product in the sequence. 
    \item (13) is to ensure a setup when period $t $ is production-free and the setup state at the end of $t - 1 $ is different from the setup state at the end of $t $
    \item (14) and (15) counts setups between periods in which the machine is not idle
\end{itemize}
\begin{gather}
    T_{ijt}\ge\alpha_{it} + \gamma_{j, t - 1}\quad\forall i\neq j, t\\
    T_{ijt}\ge\alpha_{it} + \beta{j, t - 1}\quad\forall i\neq j, t\\
    0 \le \omega_t\le 1\quad\forall t\\
    X_{it}, I_{it}, \delta_{t}\ge0\quad\forall i, t\\
    T_{ijt}, Y_{it}, \alpha_{it}, \beta_{it}, \gamma_{it}\in\{0, 1\}\quad\forall i, t\\
    Y_{i, t} + \gamma_{i, t - 1} \ge \sum_{j = 1}^{N}T_{i, j, t}\quad\forall i, t\\
    Y_{j, t} + \gamma_{j, t} \ge \sum_{j = 1}^{N}T_{i, j, t}\quad\forall j, t\\
    \beta_{it} + \alpha_{i, t + 1} - 1 \le \gamma_{it}
\end{gather}
\newpage 
\section{A simple optimization problem}
A company has one product
\begin{itemize}
    \item The sales price is 500 per unit 
    \item The cost is 200 per unit 
    \item It takes 5 hours to produce one unit
    \item There is a total production cap of 750 hours 
    \item The market will buy all units that the company sells 
\end{itemize}
How many units should be produced and sold? What will be the total profit?

\begin{gather*}
    \max(500 - 200)X\\
    \text{s.t.}\\
    5X \le 750\\
    0\le X \le\infty
\end{gather*}
\section{A slightly more complex problem}
A company has two products, 1 and 2
\begin{itemize}
    \item $P_1 = 500\quad P_2 = 400 $
    \item $C_1 = 200\quad C_2 = 150 $
    \item $T_1 = 5\quad T_2 = 3 $
    \item There is a total capacity of 750 hours ($5X_1 + 3X_2 \le 750 $)
    \item Maximum sales of the products are 120 for $P_1 $ and 200 for $P_2 $ 
    \begin{itemize}
        \item $X_1 \le 120\quad X_2 \le 200 $
    \end{itemize}
\end{itemize}
Okay so basically (without looking at the solution):
\begin{gather*}
    \max(500 - 200)X_1 + (400 - 150)X_2\\
    \text{s.t.}\\
    0 \le X_1 \le 120\\
    0 \le X_2 \le 200\\ 
    5X_1 + 3X_2 \le 750
\end{gather*}
\textbf{Set-based form}
\begin{gather*}
    \max\left(\sum_{i}(P_i - C_i)X_i\right)\\
    \text{s.t.}\\
    \sum_{i}T_iX_i\le \text{capacity}\quad\text{(here, 750)}\\
    0\le X_i \le \text{max sales}\quad\forall i \quad\text{(here, 120 and 200)}
\end{gather*}
\newpage 
\section{Demand-allocation model}

A product is produced at three factories $(A, B, \text{ and } C)$ and shipped to four markets, (1, 2, 3, 4)

Capacities at the three factories: 
\begin{center}
    \begin{tabular}{|c|c|c|}
        \hline 
        \textbf{\textit{A}} & \textbf{\textit{B}} & \textbf{\textit{C}} \\\hline 
        100 & 120 & 110 \\\hline 
    \end{tabular}
\end{center}

Demand at the four markets: 
\begin{center}
    \begin{tabular}{|c|c|c|c|}
        \hline 
        \textbf{1} & \textbf{2} & \textbf{3} & \textbf{4} \\\hline 
        55 & 60 & 35 & 50 \\\hline 
    \end{tabular}
\end{center}

Cost per unit for shipping: 
\begin{center}
    \begin{tabular}{|c|c|c|c|c|}
        \hline 
        & \textbf{1} & \textbf{2} & \textbf{3} & \textbf{4} \\\hline 
        \textbf{A} & 7 & 4 & 9 & 6 \\\hline 
        \textbf{B} & 3 & 8 & 5 & 4 \\\hline 
        \textbf{C} & 6 & 5 & 10 & 5 \\\hline 
    \end{tabular}
\end{center}

\textbf{Formulate a mathematical model to find the optimal quantities to ship from each factory to each market, minimizing total costs}

Parameters: 
\begin{itemize}
    \item $c_{i, j}$: Unit cost from factory $i $ to market $j $
    \item $d_j$: Demand at market $j $
    \item $\text{cap}_i $: Capacity at factory $i $
\end{itemize}

Decision variables: 
\begin{itemize}
    \item $X_{i, j}$: quantity shipped from factory $i $ to market $j $
\end{itemize}

\begin{gather*}
    \min\left(\sum_{i}\sum_{j}c_{i, j}X_{i, j}\right)\\
    \text{s.t.}\\
    \sum_{j}X_{i, j}\le\text{cap}_{i}\quad\forall i\in\{A, B, C\}\\
    \sum_{i}X_{i, j} = d_j\quad\forall j \in \{1, 2, 3, 4\}\\
    X_{j, i}\ge 0
\end{gather*}

\end{document}